\documentclass[../../../include/open-logic-section]{subfiles}

\begin{document}

\olfileid{cmp}{thy}{rus}
\olsection{与Russell悖论的比较}

将本节的论证与Russell悖论加以比较和对照，颇具启发性：
% Part: computability
% Chapter: computability-theory
% Section: russells-paradox

\begin{enumerate}
\item Russell悖论：令 $S = \Setabs{x}{x \notin x}$。则 $S \in S$ 当且仅当 $S \notin S$，矛盾。

  \emph{结论：}不存在这样的集合~$S$。假设“所有集合的集合”存在，会与集合论的其他公理不一致。

\item Russell悖论的一个变体：令 $F$ 为从所有函数的集合到 $\{ 0, 1 \}$ 的一个“函数”，定义为
  \[
  F(f) =
  \begin{cases}
    1 & \text{若 $f$ 在 $f$ 的定义域内，且 $f(f) = 0$} \\
    0 & \text{否则}
  \end{cases}
  \]
  类似的论证表明 $F(F) = 0$ 当且仅当 $F(F) = 1$，矛盾。

  \emph{结论：}$F$ 不是一个函数。“所有函数的集合”太大，无法作为函数的定义域。

\item 对角线论证：令 $f_0$, $f_1$, \dots 是部分可计算函数的枚举，并令 $G \colon \Nat \to
  \{ 0, 1 \}$ 定义为
  \[
  G(x) =
  \begin{cases}
    1 & \text{若 $f_x(x)\downarrow = 0$} \\
    0 & \text{否则}
  \end{cases}
  \]
  若 $G$ 是可计算的，则对某个 $k$，它是函数 $f_k$。但如此一来 $G(k) = 1$ 当且仅当 $G(k) = 0$，矛盾。

  \emph{结论：}$G$ 不可计算。注意，根据集合论的公理，$G$ 仍然是一个函数；这里并无悖论，只是一点澄清。
\end{enumerate}

关于偏函数、可计算函数、部分可计算函数等的讨论可能会令人困惑。从 $\Nat$ 到 $\Nat$ 的所有偏函数构成一个庞大的对象集。其中有些是全函数，有些是可计算的，有些既是全函数又是可计算的，还有些两者都不是。请注意，当我们说“函数”时，默认指的是全函数。因此我们有：
\begin{enumerate}
\item 可计算函数
\item 非全的部分可计算函数
\item 不可计算的全函数
\item 既不全也不可计算的偏函数
\end{enumerate}

为了理清这些概念，不妨画一个大正方形来代表所有从 $\Nat$ 到 $\Nat$ 的偏函数，然后标出两个互相重叠的区域，分别对应全函数和可计算的偏函数。这是一个很好的练习：试着在图中每一个最终形成的区域里描述出一个对象。

\end{document}